We now prove some facts about the quivers $\tqdn$, which will be useful in proving the main theorem of this section. These will be used in combination with the previous two lemmas to show that one can realise $\mathcal{Q}(\mathcal{T})$ as a slice if $\mathcal{T}$ has no interior $(d + 1)$-simplices. We say that a full subquiver $P$ of $\tqdn$ is \emph{switching-closed} if whenever $A \rightarrow A + E_{i} \rightarrow A + E_{i} + E_{j}$ is a sequence of arrows in $P$ and $A + E_{j} \in \tqdn_{0}$, then the sequence of arrows $A \rightarrow A + E_{j} \rightarrow A + E_{i} + E_{j}$ is also in $P$.

\begin{lemma}\label{lem:all_paths}
Let $P$ be a switching-closed full subquiver of $\tqdn$. If there is a path $A \leadsto B$ in $P$, all other paths $A \leadsto B$ in $\tqdn$ must also lie in $P$.
\end{lemma}
\begin{proof}
Suppose that we have a path $A \leadsto B$ in $P$. The length of any such path is $\sum_{i = 0}^{d}(b_{i} - a_{i})$. We prove the claim by induction on this quantity. The base case, where the length is $1$, follows from the fact that $P$ is a full subquiver of $\tqdn$.

For the inductive step, we assume that the claim holds for all $X$ and $Y$ with $\sum_{i = 0}^{d}(y_{i} - x_{i}) < \sum_{i = 0}^{d}(b_{i} - a_{i})$. We have that $B$ is the head of up to $(d + 1)$ arrows, namely the ones with tails $B - E_{0}, B - E_{1}, \dots, B - E_{d}$, provided these are vertices of $\tqdn$. Suppose that $B - E_{i}$ is the penultimate vertex of our path $A \leadsto B$.

Choose a vertex $B - E_{j} \in \tqdn_{0}$, where $i \neq j$. If $A \not\leqslant B - E_{j}$, then we may ignore this vertex, since there can be no paths from $A$ to $B$ through it. Hence we assume that $A \leqslant B - E_{j}$. This implies that $A \leqslant B - E_{i} - E_{j} \leqslant B - E_{i}$, so $B - E_{i} - E_{j} \in P_{0}$ by the induction hypothesis applied to $A$ and $B - E_{i}$.

Since $P$ is switching-closed, we have that $B - E_{j} \in P_{0}$, since $B - E_{i} - E_{j}, B - E_{i}, B \in P_{0}$. By the induction hypothesis, all paths $A \leadsto B - E_{j}$ lie in $P$, and hence all paths $A \leadsto B$ passing through $B - E_{j}$ lie in $P$. The result follows.
\end{proof}

By a \emph{walk} in a quiver $Q$ from $A$ to $B$ we mean a finite sequence of arrows $\beta_{1}\beta_{2}\dots\beta_{s}$ such that $\beta_{1}$ is incident at $A$, $\beta_{s}$ is incident at $B$, and $\beta_{i-1}$ and $\beta_{i}$ are incident at a common vertex for all $i \in \{2, 3, \dots, s\}$. In this case, we write $A \walk B$. That is, a path only consists of forwards arrows, but a walk may contain backwards arrows as well.

\begin{lemma}\label{lem:walk_to_path}
Let $P$ be a connected switching-closed full subquiver of $\tqdn$. If $A,B \in P_{0}$ are such that there is a path $A \leadsto B$ in $\tqdn$, then there is a path $A \leadsto B$ in $P$.
\end{lemma}
\begin{proof}
Let $A,B \in P_{0}$. Suppose that there is a path $A \leadsto B$ in $\tqdn$. There is certainly a walk $W \colon A \walk B$ in $P$, since $P$ is connected. We prove that there is also a path $A \leadsto B$ in $P$ by induction on the number of backwards arrows in this walk. The base case, in which there are zero backwards arrows in the walk, is immediate.

Hence we suppose for induction that the claim holds for walks with fewer backwards arrows than $W$. We may assume that the final arrow in $W$ is a backwards one, otherwise we may remove the final arrow and consider instead the walk $A \walk B'$, where $A \walk B' \rightarrow B$ is the original walk~$W$.

Therefore we can assume that our walk is of the form $A \walk C \leftarrow B$, where $C \in P$. By the induction hypothesis, we can replace this with a walk of the form $A \leadsto C \leftarrow B$ in $P$. Moreover, by Lemma~\ref{lem:all_paths} we have that all paths $A \leadsto C$ in $\tqdn$ lie in $P$. Since we have a path $A \leadsto B$ in $\tqdn$, we have $A \leqslant B$ and, moreover, $A \leqslant B \leqslant C$. There is therefore a path $A \leadsto B \leadsto C$ in $\tqdn$. This path is in $P$, since every path $A \leadsto C$ is in $P$, which gives the desired path $A \leadsto B$ in~$P$.
\end{proof}

These two lemmas imply the following corollary.

\begin{coro}\label{cor:switch+con->conv}
Let $P$ be a connected switching-closed full subquiver of $\tqdn$. Then $P$ is convex in $\tqdn$.
\end{coro}
\begin{proof}
Suppose that $P$ is a connected switching-closed full subquiver of $\tqdn$. Let $A, B \in P_{0}$ be such that there is a path $A \leadsto B$ in $\tqdn$. By Lemma~\ref{lem:walk_to_path}, there is a path $A \leadsto B$ in $P$. Then, by Lemma~\ref{lem:all_paths}, we have that all paths $A \leadsto B$ lie in $P$, and so $P$ is convex.
\end{proof}

This corollary is useful because it is easier to check the property of being connected and switching-closed than the property of being convex. Given a full subquiver $P$ of $\tqdn$, we write $\overline{P}$ for the smallest switching-closed subquiver containing $P$. The subquiver $\overline{P}$ is well-defined since the intersection of a set of switching-closed full subquivers is switching-closed, so $\overline{P}$ may be constructed as the intersection of all switching-closed subquivers containing $P$.
